\section{Conclusion}

We have seen that extending a derivative operation to \( \Type \)-valued containers is possible:
If the obstruction to a general derivative is that positions are not necessarily discrete,
simply restrict a derivative to the subtype of those positions that are!
This generalization does not come for free, however.
Although it is universal and retains basic properties, the chain rule is no longer in general invertible,
and neither is the \( \mu \)-rule.
In fact, a global chain rule is impossible, and at least constructively taboo when restricted to sets.

Interestingly, the status of the \( \mu \)-rule is somewhat weaker:
In case of the chain rule the contradiction arises because we are free to apply it to an arbitrary pair of containers;
in \autoref{globally-discrete-iff-strong-chain-rule} the strategy is to pick an \enquote{adversarial}
container \( G \) which forces the positions of an arbitrary container \( F \) to be discrete.
In \autoref{strong-mu-rule-iff-strong-chain-rule}, however, the constraints are tighter
---
there is one degree of freedom (the container \( F \)),
and whether the rule is strong depends on that and \( \mu F \),
which is canonically derived from \( F \).
We would like to know:
Is a strong \( \mu \)-rule a constructive taboo in the same way that a strong chain rule is?
We currently do not know of any container \( F \) for which \( \MuRule_F \) fails to be strong.
At the same time, proving that it is always strong may require more advanced proof techniques,
such as the generalized \( \WPath \)-induction of~\cite[Lemma~12]{DamatoEtAl2025FormalisingInductiveCoinductive}.

In the future we would also like to work out the behaviour of \( \Der \) on fixpoints other than \( \mu \).
In their original work~\cite{AbbottEtAl2005DataDifferentiating}, \citeauthor{AbbottEtAl2005DataDifferentiating} derive a fixpoint rule also for the greatest fixpoint, \( \nu F \).
Recently, \citeauthor{DamatoEtAl2025FormalisingInductiveCoinductive} show that \( \Type \)-valued container functors preserve both least and greatest fixpoints~\cite{DamatoEtAl2025FormalisingInductiveCoinductive}.
In either work, many of the intermediate lemmata hold for arbitrary fixpoints,
and the same is true for a number of our results:
The construction of the lax chain rule in \autoref{lax-chain-rule}, for example, only depends on the property of \( \In_F : \Cart{\Subst{F}{\mu F}}{\mu F} \) being \emph{some}
\( \Subst{F}{\Blank} \)-fixpoint, not necessarily the least fixpoint.
We could thus give for any \( \varphi : \CartIso{\Subst{F}{\varphi F}}{\varphi F} \) a \enquote{\( \varphi \)-rule}
as an embedding \( \operatorname{\Op{\varphi-rule}} : \Cart{\mu F'}{\Der(\varphi F)} \) in which
\(
  F' \DefEq
    { \Wk{\Der_0{F}[ {\varphi F} ]} }
      \CPlus
    (
      \Wk{\Der_1{F}[ {\varphi F} ]}
        \CTimes
      \Proj{1}
    )
\).

In another direction, we are curious whether we can describe more than just one-hole contexts.
\Citeauthor{AbbottEtAl2003DerivativesContainers} define \emph{linear exponentials} \cite[{Def.~4.1}]{AbbottEtAl2003DerivativesContainers}:
for containers \( F \) and \( H \), the linear exponential of \( [H, F] \) exists if there is a universal morphism from \( \Blank \CTimes H \) to \( F \).
They show that \( \Der{F} = [\Id, F] \) if \( F \) is discrete.
The adjunction of \autoref{derivative-adjunction-sets} is such a universal morphism, natural in all \( F \),
hence \( \Der = [\Id, \Blank] \) as functors.
Iterating the adjunction, we can describe \( n \)-hole contexts as linear exponentials \( \Der^n = [\CYo\Fin{n}, \Blank] \).
We conjecture that \( [\CYo(H), \Blank] \) exists for arbitrary \( H : \Type \):
for any embedding \( f : H \hookrightarrow P_s \) whose fibers are decidable,
the complement \( P_s -_f H \DefEq \sum_{p : P_s} \neg \Fiber_f(p) \) satisfies, among other things, \( (P_s -_f H) + H \simeq P_s \).
This indicates that a type of contexts with \enquote{\( H \)-shaped} holes may exist, independently of the truncation level of \( H \).
